% Generated draft: requires mathematical and source review.
\documentclass[11pt]{article}
\usepackage[margin=25mm]{geometry}
\usepackage{fontspec}
\setmainfont{Latin Modern Roman}
\usepackage{amsmath,amssymb,mathtools,unicode-math}
\setmathfont{Latin Modern Math}
\usepackage{enumitem,needspace,xurl,hyperref}
\setlength{\parindent}{0pt}
\setlength{\parskip}{5pt}
\setlength{\emergencystretch}{3em}
\newcommand{\N}{\mathbb{N}}
\newcommand{\Z}{\mathbb{Z}}
\newcommand{\Q}{\mathbb{Q}}
\newcommand{\R}{\mathbb{R}}
\newcommand{\C}{\mathbb{C}}
\newcommand{\F}{\mathbb{F}}
\newcommand{\A}{\mathbb{A}}
\newcommand{\PP}{\mathbb{P}}
\newcommand{\E}{\mathbb{E}}
\newcommand{\eps}{\varepsilon}
\newcommand{\dd}{\,\mathrm{d}}
\begin{document}
\section*{UnsolvedMath ID: 30006411}
\textbf{Source problem code:} catalogue entry
\subsection{Definitions and mathematical statement}
Let \(n\geq 3\) be an integer and let \(\delta\) denote the Euclidean metric on \(\R^n\). An \(L^\infty\) Riemannian metric uniformly bilipschitz to \(\delta\) is a measurable symmetric tensor \(g\) for which there is a constant \(C\geq 1\) such that
\[
 C^{-1}|v|^2\leq g_x(v,v)\leq C|v|^2
 \qquad\text{for almost every }x\in\R^n\text{ and every }v\in\R^n.
\]
The source considers such metrics admitting a representative smooth away from a singular set \(S\subseteq\R^n\).

For \(r>0\), write
\[
 S_r=\{x\in\R^n:\operatorname{dist}_{\delta}(x,S)<r\}.
\]
A usual upper \((n-2)\)-dimensional Minkowski-content quantity is
\[
 \limsup_{r\downarrow0}\frac{\mathcal L^n(S_r)}{r^2},
\]
where \(\mathcal L^n\) denotes Euclidean Lebesgue measure. The source does not specify whether its Minkowski content uses this global upper-content convention, an actual limit, or a local convention. Accordingly, let \(\mathcal M^{n-2}(S)\) denote the content intended by the source, pending verification of that definition.

Let \(\mathcal P_1(g,S)\) stand for the exact conclusion of the source's ``Theorem 1,'' described in the cleaned record as nonnegative-scalar-curvature smoothing. This is a placeholder: the theorem's conclusion, approximation topology, curvature requirements, and any additional hypotheses are not included in the supplied record.

With those unresolved meanings retained explicitly, the questions are whether there exist \(S\subseteq\R^n\) and an \(L^\infty\) metric \(g\), uniformly bilipschitz to \(\delta\) and smooth away from \(S\), satisfying respectively
\[
 \mathcal M^{n-2}(S)<\infty
 \quad\text{and}\quad \neg\mathcal P_1(g,S),
\]
and
\[
 \mathcal M^{n-2}(S)=0
 \quad\text{and}\quad \neg\mathcal P_1(g,S).
\]
The source also asks whether the conclusion of Theorem 1 still holds in the zero-content case. Any counterexample must satisfy the other hypotheses of that theorem; those hypotheses cannot be reconstructed from the supplied text. In particular, the cited work's title refers to nonnegative scalar curvature outside the singular set, but the displayed source question does not explicitly impose \(R_g\geq0\) on \(\R^n\setminus S\). Whether that hypothesis is inherited must be checked. These are therefore conditional formulations of the source questions, not a complete independently specified smoothing problem.
\subsection{Short English statement}
For dimensions \(n\geq3\), the source asks whether a bounded metric on \(\R^n\), uniformly comparable to the Euclidean metric and smooth outside a singular set of finite \((n-2)\)-dimensional Minkowski content, can violate a cited scalar-curvature smoothing theorem. It asks the same question for zero content. The missing theorem and content convention must be recovered before a fully precise statement is possible.
\subsection{Sources}
\begin{itemize}
\item Local UnsolvedMath registry: \texttt{data/unsolvedmath/problems.json}, with \texttt{catalog/problems.json} as the additional-record catalog; permanent UnsolvedMath ID \(30006411\), title ``Scalar-Curvature Smoothing at Critical-Codimension Singular Sets,'' exact source problem code \texttt{catalogue entry}. This draft uses the supplied complete record; no registry file or external source was independently inspected.
\item Supplied report citation: \emph{Differentialgeometrie im Großen (2026)}, Oberwolfach Reports, DOI 10.4171/owr/2025/30, pp.\ 1591--1646. \url{https://doi.org/10.4171/owr/2025/30}. The literature-source label is ``Oberwolfach Report 30/2025.'' The supplied proposed year is 2025; the citation's 2026 designation is retained without resolving the discrepancy.
\item Supplied literature reference: \emph{Smoothing L infinity Riemannian metrics with nonnegative scalar curvature outside of a singular set}. \url{https://arxiv.org/abs/2406.04564}. The original statement associates this title with Paula Burkhardt-Guim. No proposer is explicitly identified, and the text of its Theorem 1 is not supplied.
\item Literature assessment checked at \texttt{2026-08-22}: supplied dated curation, not independently verified. The supplied status is ``open.'' The supplied assessment says that the cited smoothing theorem treats singular sets of strictly smaller effective dimension and that the critical finite \((n-2)\)-dimensional Minkowski-content case and possible counterexamples remain unresolved. The background explicitly calls this dated literature triage derived from the supplied OWR curation and requests independent verification of time-sensitive status and cited claims.
\end{itemize}
\end{document}
